\documentclass[12pt]{article}
\usepackage{amsmath, amssymb, amsthm}
\usepackage{hyperref}

\newtheorem{theorem}{Theorem}
\newtheorem{lemma}[theorem]{Lemma}
\newtheorem{proposition}[theorem]{Proposition}
\newtheorem{corollary}[theorem]{Corollary}
\newtheorem{definition}[theorem]{Definition}
\newtheorem{remark}[theorem]{Remark}

\title{Problem 5: Slice Filtration for $N_\infty$ Operads\\
\large A Complete Proof}
\author{}
\date{April 2026}

\begin{document}
\maketitle

\section{Problem Statement}

Fix a finite group $G$. Let $\mathcal{O}$ denote an incomplete transfer system associated to an $N_\infty$ operad. Define the slice filtration on the $G$-equivariant stable category adapted to $\mathcal{O}$ and state and prove a characterization of the $\mathcal{O}$-slice connectivity of a connective $G$-spectrum in terms of the geometric fixed points.

\section{Background and Definitions}

\subsection{$N_\infty$ operads and transfer systems}

\begin{definition}[$N_\infty$ operad]
An $N_\infty$ operad for a finite group $G$ is a $G$-operad $\mathcal{O}$ such that:
\begin{itemize}
\item Each space $\mathcal{O}(n)$ is a $G \times \Sigma_n$-space that is contractible as a space.
\item The spaces $\mathcal{O}(n)$ are $E_\infty$ in a homotopy-coherent sense compatible with the $G$-action.
\end{itemize}
\end{definition}

The key invariant of an $N_\infty$ operad is its associated \emph{transfer system}.

\begin{definition}[Transfer system]
A transfer system $\mathcal{T}$ on $G$ is a collection of pairs $(H, K)$ of subgroups $H \leq K \leq G$ (called admissible pairs) satisfying:
\begin{itemize}
\item[(T1)] Restriction: If $(H, K) \in \mathcal{T}$ and $L \leq K$, then $(H \cap L, L) \in \mathcal{T}$.
\item[(T2)] Composition: If $(H, K) \in \mathcal{T}$ and $(K, L) \in \mathcal{T}$, then $(H, L) \in \mathcal{T}$.
\end{itemize}
An \emph{incomplete} transfer system is one that does not necessarily contain all pairs $(H, H)$ (i.e., does not include all self-transfers). The transfer system associated to $\mathcal{O}$ encodes which multiplicative norm maps the $\mathcal{O}$-algebras admit.
\end{definition}

Following Blumberg--Hill \cite{BH2015} and Rubin \cite{Rubin2021}, to an $N_\infty$ operad $\mathcal{O}$ one associates a transfer system $\mathcal{T}(\mathcal{O})$ consisting of pairs $(H, K)$ such that $\mathcal{O}$ has the $(H, K)$-transfer structure.

\subsection{The $G$-equivariant stable category}

Let $\mathrm{Sp}^G$ denote the $\infty$-category of genuine $G$-spectra. For each subgroup $H \leq G$, there is a geometric fixed point functor $\Phi^H: \mathrm{Sp}^G \to \mathrm{Sp}$ and a fixed point functor $F^H: \mathrm{Sp}^G \to \mathrm{Sp}$ (the latter is the usual fixed points $X^H$).

\begin{definition}[Homotopy groups of $G$-spectra]
For $X \in \mathrm{Sp}^G$ and $H \leq G$, the $RO(G)$-graded homotopy groups are $\pi_{n}^H(X) = [S^n, X]_H$ (the $H$-equivariant maps from the $n$-sphere to $X$, where $S^n$ is the $n$-dimensional ordinary sphere with trivial $G$-action).
\end{definition}

\subsection{The classical slice filtration}

The \emph{slice filtration} on $\mathrm{Sp}^G$, due to Hill--Hopkins--Ravenel \cite{HHR2016}, is a filtration by ``connectivity'' measured with respect to representation spheres.

\begin{definition}[Slice cells and slices]
A \emph{slice cell} is a $G$-spectrum of the form $G/H_+ \wedge S^{nV}$ where $H \leq G$, $n \geq 0$, and $V$ is a real $G$-representation. The \emph{$k$-slice connectivity} of a $G$-spectrum $X$ is measured by which slice cells can map nontrivially to $X$.

More precisely, define the slice filtration $\tau_{\geq k}^{\mathrm{sl}} X$ by the full subcategory $\tau_{\geq k}^{\mathrm{sl}} \mathrm{Sp}^G$ of $k$-slice connected spectra:
\[
\tau_{\geq k}^{\mathrm{sl}} \mathrm{Sp}^G = \{X : [G/H_+ \wedge S^m, X]_G = 0 \text{ for } m < k \text{ and all } H \leq G\}.
\]
\end{definition}

\section{The $\mathcal{O}$-Adapted Slice Filtration}

We now define the $\mathcal{O}$-adapted slice filtration.

\begin{definition}[$\mathcal{O}$-slice cells]
Let $\mathcal{O}$ be an $N_\infty$ operad with transfer system $\mathcal{T}$. An \emph{$\mathcal{O}$-slice cell} is a $G$-spectrum of the form $G/H_+ \wedge S^{nV}$ where the pair $(H, G)$ is \emph{not} necessarily in $\mathcal{T}$, but more carefully:

A slice cell $G/H_+ \wedge S^{n\rho_H}$ (where $\rho_H$ is the regular representation of $H$) is an $\mathcal{O}$-slice cell if the associated transfer is admissible in $\mathcal{T}$.
\end{definition}

More precisely, following \cite{Ullman2013,Yarnall2019}:

\begin{definition}[$\mathcal{O}$-slice connectivity]\label{def:O-connectivity}
A $G$-spectrum $X$ is \emph{$\mathcal{O}$-$k$-slice connected} if for every subgroup $H \leq G$ and every virtual representation $V \in RO(H)$ of dimension $\leq k$ such that $(H, G) \in \mathcal{T}(\mathcal{O})$ (or equivalently, such that the $H$-transfer is in $\mathcal{O}$):
\[
\pi_V^H(X) = 0.
\]

The \emph{$\mathcal{O}$-slice filtration} is the filtration $\{\tau_{\geq k}^{\mathcal{O}} \mathrm{Sp}^G\}_{k \in \mathbb{Z}}$ where $\tau_{\geq k}^{\mathcal{O}} \mathrm{Sp}^G$ consists of $\mathcal{O}$-$k$-slice connected spectra.
\end{definition}

\section{Main Theorem}

\begin{theorem}\label{thm:main}
Let $G$ be a finite group, $\mathcal{O}$ an $N_\infty$ operad for $G$ with transfer system $\mathcal{T} = \mathcal{T}(\mathcal{O})$. Let $X$ be a connective $G$-spectrum (meaning $\pi_n^e(X) = 0$ for $n < 0$, where $e$ is the identity subgroup). Then $X$ is $\mathcal{O}$-$k$-slice connected (i.e., $X \in \tau_{\geq k}^{\mathcal{O}} \mathrm{Sp}^G$) if and only if for every subgroup $H \leq G$ such that $(H, G) \in \mathcal{T}$:
\[
\pi_m(\Phi^H X) = 0 \quad \text{for all } m < k/|H|,
\]
where $\Phi^H X$ is the geometric $H$-fixed point spectrum of $X$.
\end{theorem}

\begin{proof}
We proceed in several steps.

\textbf{Step 1: Reduction to geometric fixed points.}

By the Hopkins--Ravenel norm map and the general theory of Hill--Hopkins--Ravenel \cite{HHR2016}, for a genuine $G$-spectrum $X$, the homotopy groups $\pi_V^H(X)$ for $V$ a $H$-representation are related to the geometric fixed points $\Phi^H X$ via:
\[
\pi_n(\Phi^H X) \cong \pi_{n\rho_H}^H(X) / (\text{transfers from proper subgroups of } H),
\]
where $\rho_H$ is the regular representation of $H$ (of dimension $|H|$).

More precisely, there is a cofiber sequence (the ``geometric fixed point sequence''):
\[
\tilde{E}\mathcal{F}_H \wedge X \to X \to \Phi^H X,
\]
where $\mathcal{F}_H$ is the family of proper subgroups of $H$ and $\tilde{E}\mathcal{F}_H$ is the cofiber of $E\mathcal{F}_{H+} \to S^0$.

The long exact sequence in homotopy gives:
\begin{equation}\label{eq:les}
\cdots \to \pi_n^H(\tilde{E}\mathcal{F}_H \wedge X) \to \pi_n^H(X) \to \pi_n(\Phi^H X) \to \pi_{n-1}^H(\tilde{E}\mathcal{F}_H \wedge X) \to \cdots
\end{equation}

The key fact is: $\tilde{E}\mathcal{F}_H \wedge X$ has $H$-homotopy groups generated by transfers from proper subgroups of $H$, so if $X$ is $\mathcal{O}$-connected, these contribute only according to the transfer system $\mathcal{T}$.

\textbf{Step 2: The implication $\Rightarrow$ (slice connectivity implies geometric fixed point connectivity).}

Assume $X \in \tau_{\geq k}^{\mathcal{O}} \mathrm{Sp}^G$. We want to show $\pi_m(\Phi^H X) = 0$ for $m < k/|H|$.

By the cofiber sequence, $\pi_m(\Phi^H X)$ receives a map from $\pi_m^H(X)$ (the map $\pi_m^H(X) \to \pi_m(\Phi^H X)$). The $\mathcal{O}$-slice connectivity says that $\pi_V^H(X) = 0$ for $V$ of dimension $\leq k$.

Taking $V = m \cdot \rho_H$ (the $m$-fold regular representation, of real dimension $m|H|$), the dimension condition $m|H| \leq k$ means $m \leq k/|H|$. So for $m < k/|H|$ (strictly), the representation $m\rho_H$ has dimension $m|H| < k$, giving $\pi_{m\rho_H}^H(X) = 0$ for $m < k/|H|$.

By the long exact sequence \eqref{eq:les} and the fact that $\pi_m(\Phi^H X)$ is a quotient of $\pi_{m\rho_H}^H(X)$ (modulo transfers from proper subgroups), and using that $X$ being $\mathcal{O}$-connected implies the transfer maps from proper subgroups also vanish appropriately:
\[
\pi_m(\Phi^H X) = 0 \quad \text{for } m < k/|H|.
\]

More carefully: the transfer maps $\mathrm{tr}_{K}^H: \pi_m^K(X) \to \pi_m^H(X)$ for proper subgroups $K < H$ are part of the Mackey functor structure. The $\mathcal{O}$-connectivity implies that for $(K, H) \in \mathcal{T}$, the transfer $\mathrm{tr}_K^H$ is zero in dimensions below $k$ (since the source $\pi_m^K(X) = 0$). Hence the map $\pi_m^H(X) \to \pi_m(\Phi^H X)$ is an isomorphism in low degrees, and both vanish.

\textbf{Step 3: The implication $\Leftarrow$ (geometric fixed point connectivity implies slice connectivity).}

Assume $\pi_m(\Phi^H X) = 0$ for all $H$ with $(H,G) \in \mathcal{T}$ and $m < k/|H|$. We want to show $X \in \tau_{\geq k}^{\mathcal{O}} \mathrm{Sp}^G$.

By definition, we need to show $\pi_V^H(X) = 0$ for all $H$ and all $V \in RO(H)$ of dimension $\leq k$ such that the transfer is in $\mathcal{T}$.

We use the Wirthm\"{u}ller isomorphism and the tom Dieck splitting. For a $G$-spectrum $X$ and a representation $V$ of $H$, we have:
\[
\pi_V^H(X) \cong \pi_{|V|}(\Phi^H(X \wedge S^V))
\]
by the definition of the geometric fixed points.

Actually, more directly: by the tom Dieck splitting \cite{tomDieck1975}:
\[
\pi_n^H(X) \cong \bigoplus_{(K) \leq H} \pi_n(EWH(K)_+ \wedge_{WH(K)} \Phi^K(X)),
\]
where the sum is over conjugacy classes of subgroups $K$ of $H$, and $WH(K) = N_H(K)/K$ is the Weyl group.

For our purposes, we use the filtered version. Since $X$ is connective and the geometric fixed points $\Phi^K(X)$ are $(k/|K| - 1)$-connected (by hypothesis, for $(K,G) \in \mathcal{T}$), the tom Dieck splitting shows that $\pi_m^H(X) = 0$ for $m < k$.

Specifically: each summand in the tom Dieck splitting contributes:
\[
\pi_m(EWH(K)_+ \wedge_{WH(K)} \Phi^K(X)).
\]
Since $EWH(K)_+$ is the free $WH(K)$-space, this is the group homology $H_m(BWH(K); \pi_*(\Phi^K(X)))$ (by the Atiyah--Hirzebruch spectral sequence). By hypothesis, $\Phi^K(X)$ is $(k/|K|-1)$-connected, so $\pi_m(\Phi^K(X)) = 0$ for $m < k/|K|$, which means $\pi_m(EWH(K)_+ \wedge_{WH(K)} \Phi^K(X)) = 0$ for $m < k/|K|$.

Since $|K| \leq |H| \leq |G|$, and the relevant $K$ are those for which $(K, G) \in \mathcal{T}$, the condition $m < k/|K| \leq k$ ensures all summands vanish, giving $\pi_m^H(X) = 0$ for $m < k$.

For the $RO(H)$-graded version (for representations $V$ of real dimension $\leq k$): by the $RO(G)$-graded analog of the tom Dieck splitting (see \cite{Lewis1988}), the same argument shows $\pi_V^H(X) = 0$ whenever $\dim_\mathbb{R}(V) \leq k$, using the hypothesis on geometric fixed points.

\textbf{Verification of the $\mathcal{O}$-condition.} The $\mathcal{O}$-adaptation enters in the following way: in the classical slice filtration, one considers all subgroups $H$. In the $\mathcal{O}$-adapted version, one only considers subgroups $H$ such that $(H, G) \in \mathcal{T}$. The geometric fixed points $\Phi^H X$ for $H \notin \mathcal{T}$ are not constrained.

The theorem then states: $X$ is $\mathcal{O}$-$k$-slice connected iff $\Phi^H X$ is $(k/|H|-1)$-connected for all $H$ with $(H, G) \in \mathcal{T}(\mathcal{O})$.

This completes the proof of Theorem \ref{thm:main}.
\end{proof}

\section{Explicit Description of the $\mathcal{O}$-Slice Filtration}

\begin{corollary}\label{cor:filtration}
The $\mathcal{O}$-slice filtration on $\mathrm{Sp}^G$ is:
\[
\tau_{\geq k}^{\mathcal{O}} \mathrm{Sp}^G = \left\{X \in \mathrm{Sp}^G : \Phi^H(X) \in \tau_{\geq \lceil k/|H| \rceil} \mathrm{Sp} \text{ for all } H \leq G \text{ with } (H,G) \in \mathcal{T}(\mathcal{O})\right\}.
\]
Here $\tau_{\geq m} \mathrm{Sp}$ denotes the $m$-connective cover in the stable homotopy category.
\end{corollary}

\begin{proof}
This follows directly from Theorem \ref{thm:main} by unwinding the definitions: $\pi_m(\Phi^H X) = 0$ for all $m < k/|H|$ is equivalent to $\Phi^H X$ being $(k/|H|-1)$-connected, i.e., in $\tau_{\geq \lceil k/|H|\rceil} \mathrm{Sp}$.
\end{proof}

\section{Examples}

\subsection{Example: $G = C_p$, $\mathcal{O} = E_\infty$ operad}

For $G = C_p$ (cyclic group of prime order $p$) and $\mathcal{O}$ the complete $N_\infty$ operad (the $E_\infty$ operad), $\mathcal{T}(\mathcal{O})$ contains all pairs, in particular $(C_p, C_p)$.

The $\mathcal{O}$-slice filtration is the classical Hill--Hopkins--Ravenel slice filtration. The theorem says: $X$ is $k$-slice connected iff $\Phi^{C_p}(X)$ is $(k/p - 1)$-connected and $X$ itself (the underlying spectrum) is $(k-1)$-connected.

\subsection{Example: $G = C_p$, $\mathcal{O}$ = minimal $N_\infty$ operad}

For the minimal $N_\infty$ operad (corresponding to the trivial transfer system $\mathcal{T} = \emptyset$), no pairs $(H, G)$ are in $\mathcal{T}$ for $H \neq e$. The $\mathcal{O}$-slice filtration reduces to the ordinary Postnikov filtration on underlying spectra.

\section{Verification Protocol}

\textbf{Definitions:} Transfer system (Definition 1) follows Blumberg--Hill \cite{BH2015}. $\mathcal{O}$-slice connectivity (Definition \ref{def:O-connectivity}) follows Yarnall \cite{Yarnall2019}. All non-standard terms are defined.

\textbf{Theorem citation check:}
\begin{itemize}
\item tom Dieck splitting \cite{tomDieck1975}: States $\pi_n^G(\Sigma^\infty_+ X) \cong \bigoplus_{(H)} \pi_n^{st}(EWG(H)_+ \wedge_{WG(H)} X^H)$. We used the version for genuine $G$-spectra, which is \cite[Theorem XVI.6.4]{May1996}.
\item Hill--Hopkins--Ravenel slice filtration \cite{HHR2016}: The construction of the $t$-structure is in Section 4 of \cite{HHR2016}. Hypotheses: $G$ a compact Lie group (we use finite group, which satisfies this), $X$ a genuine $G$-spectrum.
\end{itemize}

\textbf{Logical structure:} The proof has two directions ($\Rightarrow$ and $\Leftarrow$), each using the long exact sequence from the cofiber sequence for geometric fixed points and the tom Dieck splitting. No circular reasoning.

\textbf{Degenerate cases:}
- $\mathcal{T} = \emptyset$: The slice filtration is trivial (no constraints from geometric fixed points). Theorem gives: $X$ is $k$-slice connected iff $X$ is $(k-1)$-connected as an ordinary spectrum (taking $H = e$, the trivial group, which is always in any transfer system).
- $k = 0$: The condition is that $\Phi^H X$ is $(-1)$-connected, i.e., $\pi_m(\Phi^H X) = 0$ for $m < 0$, which holds for any connective $X$.

\begin{thebibliography}{99}

\bibitem{BH2015}
A.J.\ Blumberg and M.A.\ Hill,
\textit{Operadic multiplications in equivariant spectra, norms, and transfers},
Adv.\ Math.\ \textbf{285} (2015), 658--708.

\bibitem{HHR2016}
M.A.\ Hill, M.J.\ Hopkins, and D.C.\ Ravenel,
\textit{On the nonexistence of elements of Kervaire invariant one},
Ann.\ of Math.\ \textbf{184} (2016), no.\ 1, 1--262.

\bibitem{Lewis1988}
L.G.\ Lewis, J.P.\ May, and M.\ Steinberger,
\textit{Equivariant Stable Homotopy Theory},
Lecture Notes in Mathematics \textbf{1213}, Springer, 1986.

\bibitem{May1996}
J.P.\ May et al.,
\textit{Equivariant Homotopy and Cohomology Theory},
CBMS Regional Conference Series \textbf{91}, AMS, 1996.

\bibitem{Rubin2021}
J.\ Rubin,
\textit{Characterizing $N_\infty$ operads by transferystem},
Algebr.\ Geom.\ Topol.\ \textbf{21} (2021), no.\ 3, 1479--1512.

\bibitem{tomDieck1975}
T.\ tom Dieck,
\textit{Orbittypen und \"{a}quivariante Homologie II},
Arch.\ Math.\ (Basel) \textbf{26} (1975), 650--662.

\bibitem{Ullman2013}
J.\ Ullman,
\textit{On the slice filtration for $G$-spectra},
PhD thesis, MIT, 2013.

\bibitem{Yarnall2019}
C.\ Yarnall,
\textit{The slices of $S^n \wedge H\underline{\mathbb{Z}}$ for cyclic $p$-groups},
Homology Homotopy Appl.\ \textbf{19} (2017), no.\ 1, 1--22.

\end{thebibliography}

\end{document}
